\documentclass[12pt]{article}
\usepackage[a4paper,margin=1in]{geometry}
\usepackage{setspace}
\setstretch{1.15}

\begin{document}

\begin{center}
    \Large \textbf{Nonlinear Correlation and Stochastic Transport in Plasma Turbulence} \\[4pt]
    \normalsize \emph{In satisfaction of the requirements for Math 402} \\
    \normalsize \textbf{Supervisor:} Dr.\ Chijin Xiao (\emph{Physics and Engineering Physics}) \hspace{1cm} 
\end{center}

\section*{Overview}
This project explores a nonlinear correlation method for studying turbulent transport in magnetically confined plasmas.
Originally proposed by Ding \textit{et al.} (Phys.\ Rev.\ Lett.\ 79, 2458, 1997), the techniques analyze conditional cross-correlations between two time-resolved data series from measurement at two adjacent locations to study the direction of energy transfer between two nonlinearly coupled oscillators. Later studies (Bsharat \textit{et al.}, 2024; van Milligen \textit{et al.}, 2019) have shown that this approach can be linked to ideas from information theory and stochastic transport models such as continuous-time random walks. The aim of this work is to better understand the mathematical and physical basis of this method and how it can be applied to real plasma data.
\section*{Aims}
\begin{enumerate}
    \item Reproduce and test the nonlinear correlation algorithm using STOR-M Langmuir probe data and synthetic coupled oscillator models.
    \item Explore links between nonlinear correlation, transfer entropy, and stochastic transport models such as continuous-time random walks. 
    \item Develop a minimal Python implementation of the algorithm based on legacy MATLAB codes. In addition, if possible, port all relevant MATLAB code to Python.
\end{enumerate}

\section*{Method and Timeline}
\begin{itemize}
	\item Phase 1 (Nov--Dec): Validate the algorithm numerically and with existing STOR-M datasets.  
	\item Phase 2 (Jan--Mar): Compute nonlinear directionality $S_{xy}$ from the data from various experiments and compare it with the standard linear cross-correlation. Check if the direction or strength of correlation changes over different time windows.	
	\item Phase 3 (Mar--Apr): Interpret results in a stochastic framework and compare with information-theoretic analyses.
\end{itemize}


\section*{Expected Outcome}
The study should clarify how nonlinear correlation measures relate to stochastic models of transport and provide a simple computational tool for analyzing fluctuations in fusion plasmas.

\section*{Preliminary References}
\begin{itemize}
    \item W.\ X.\ Ding \textit{et al.}, \textit{Phys.\ Rev.\ Lett.} \textbf{79}, 2458 (1997).
    \item H.\ Bsharat \textit{et al.}, \textit{Radiation Effects \& Defects in Solids} \textbf{179}, 1527 (2024).
    \item B.\ van Milligen \textit{et al.}, \textit{Entropy} \textbf{21}, 148 (2019).
\end{itemize}

\end{document}
